\chapter{Definição do Problema}
As Atividades Complementares (AC) constituem um requisito obrigatório para a integralização curricular no curso de Ciência da Computação da UESC. Seu cumprimento envolve a participação do discente em eventos científicos, estágios, projetos de extensão, monitorias e outras vivências formativas. Para que essas horas sejam reconhecidas oficialmente pelo Colegiado do Curso (COLCIC), o aluno deve organizar e submeter o Barema: um documento unificado que reúne os comprovantes das atividades realizadas, categoriza-os segundo as normas vigentes e totaliza a carga horária acumulada por categoria.

O processo de elaboração desse documento, contudo, é inteiramente manual e descentralizado. Não existe ferramenta oficial fornecida pela instituição para auxiliar nessa tarefa. O discente precisa reunir individualmente cada certificado em formato PDF, organizá-los em um único arquivo, preencher manualmente uma tabela de indexação — indicando em qual página do documento unificado cada comprovante se encontra — e calcular as cargas horárias respeitando os tetos máximos permitidos por categoria, conforme as regras específicas do curso de bacharelado em Ciência da Computação da UESC. Cada uma dessas etapas é executada de forma independente, sem qualquer validação automatizada, e qualquer alteração na ordem ou no conjunto de documentos invalida toda a indexação previamente realizada.

%DEGAS - O FINALZINHO DO PARÁGRAFO ACIMA... QUALQUER ALTERAÇÃO TRAZ UM RISCO, NÃO?

Esse cenário resulta em três problemas concretos e recorrentes. O primeiro é a complexidade regulatória: os alunos encontram dificuldade em aplicar corretamente as regras de carga horária definidas pelo COLCIC — como o limite de horas permitido por categoria de atividade —, o que leva a cálculos errôneos e à inclusão indevida de horas que excedem os tetos regulamentares. O segundo é a propensão a erros de indexação: o preenchimento manual da coluna de páginas é suscetível a falhas simples, porém de alto impacto — um único certificado inserido ou removido do conjunto invalida todos os números de página subsequentes, obrigando o discente a reiniciar o processo de edição, reordenação e reescrita da tabela integralmente. Relatos colhidos junto a discentes evidenciam casos em que o documento precisou ser refeito três vezes consecutivas por conta exclusivamente desse tipo de inconsistência. O terceiro problema é o retrabalho administrativo gerado na coordenação: o envio de baremas incorretos — seja por erros de cálculo, por indexação equivocada ou por separação indevida entre o barema e os documentos comprobatórios — obriga a coordenação a realizar múltiplas rodadas de revisão com cada aluno, atrasando o aproveitamento das atividades e, consequentemente, o processo de colação de grau.

É importante destacar que as soluções genéricas de manipulação de PDF disponíveis no mercado, como o iLovePDF \cite{}, não resolvem esse problema. Essas ferramentas atuam como simples agregadores de arquivos: permitem unir PDFs, mas não possuem qualquer conhecimento das regras específicas do COLCIC, não calculam cargas horárias, não aplicam tetos regulamentares e não geram a tabela de indexação automaticamente. A responsabilidade sobre toda a lógica de negócio permanece, portanto, integralmente com o aluno. O problema não é a ausência de uma ferramenta de junção de arquivos, é a ausência de um sistema que compreenda as regras do processo e as aplique de forma integrada à manipulação dos documentos.


